\documentclass[11pt]{article}

% "final" gives a non-anonymous version for the course
\usepackage[final]{acl}

\usepackage{times}
\usepackage{latexsym}
\usepackage[T1]{fontenc}
\usepackage[utf8]{inputenc}
\usepackage{microtype}
\usepackage{inconsolata}
\usepackage{booktabs}
\usepackage{amsmath}

% black links only
\hypersetup{colorlinks=true, citecolor=black, linkcolor=black, urlcolor=black}

\title{Project Proposal: Reproducing OWL and Testing Whether\\
Outlier Weighed Layerwise Sparsity Still Helps Small Modern LLMs}

\author{Tuan Khang Phan \and Xuan Gia Han Nguyen \\
  University of South Florida}

\begin{document}
\maketitle

\section{Introduction}

We will reproduce ``Outlier Weighed Layerwise Sparsity (OWL): A Missing Secret Sauce for Pruning LLMs to High Sparsity'' \citep{yin2024owl}. The paper studies one-shot unstructured pruning of LLMs with no retraining. Strong methods such as SparseGPT \citep{frantar2023sparsegpt} and Wanda \citep{sun2024wanda} prune every layer to the same sparsity. OWL questions this choice and links pruning quality to outlier features, a few hidden dimensions with very large activations \citep{kovaleva2021bertbusters,dettmers2022llmint8}.

OWL scores each weight with the Wanda metric $A_{ij} = \lVert \mathbf{X}_j \rVert_2 \cdot |W_{ij}|$. For each Transformer block $l$ it computes the Layerwise Outlier Distribution (LOD) $D^l$, the fraction of weights whose score is over $M$ times the block mean. Blocks with more outliers are pruned less: $S_l \propto 1 - D^l$, rescaled into a band of width $2\lambda$ whose mean is the target $S$. This adds almost no cost on top of Wanda or SparseGPT.

We target four claims. (C1) OWL improves both Wanda and SparseGPT at high sparsity. At 70\% on LLaMA-7B, WikiText perplexity drops from 85.77 to 24.55 with Wanda and from 26.30 to 19.49 with SparseGPT. (C2) The gain grows as the model gets smaller: 61.22 points at 7B but only 6.62 at 65B. (C3) Zero-shot accuracy also improves on seven tasks. (C4) The LOD of dense LLaMA models is far from uniform and roughly U-shaped, and good pruning methods keep outliers.

\section{Reproducing Main Results}

We reproduce C1, C3, C4 and the 7B and 13B part of C2. We skip the 30B and 65B models, the CPU speedup study, LoRA fine-tuning and the appendix studies.

\paragraph{Models.} Our main model is LLaMA-7B \citep{touvron2023llama}. The checkpoint in the official README was removed, so we use \texttt{huggyllama/llama-7b}. We also run OPT-6.7B \citep{zhang2022opt} and LLaMA-13B, with the paper's hyperparameters ($M{=}5$, $\lambda{=}0.08$ for 7B, $M{=}7$ for 13B, $M{=}10$ for OPT).

\paragraph{Data.} Calibration uses 128 random 2048-token sequences from the first C4 shard \citep{raffel2020c4}. Perplexity is measured on WikiText-2 \citep{merity2017wikitext} with 2048-token windows. The paper calls this split ``validation'', but the code uses the test split, so we follow the code. Zero-shot accuracy covers BoolQ \citep{clark2019boolq}, RTE \citep{wang2018glue}, HellaSwag \citep{zellers2019hellaswag}, WinoGrande \citep{sakaguchi2021winogrande}, ARC-e, ARC-c \citep{clark2018arc} and OpenBookQA \citep{mihaylov2018obqa}, run with lm-eval-harness \citep{gao2024lmeval}.

\paragraph{Experiments.} (R1) Table~3: magnitude pruning, plus Wanda and SparseGPT each with uniform and OWL sparsity, at 70\% on all three models. (R2) Figure~2 and Table~8: a 40\% to 80\% sweep on LLaMA-7B with Wanda, comparing OWL with uniform, global, ER, ER-plus and OWL-inverse allocations. (R3) Table~4: zero-shot accuracy of the dense and 70\% sparse LLaMA-7B models. (R4) Table~1 and Figure~1: the LOD of each dense model and how each pruning method changes it.

\paragraph{Known risks.} The code's zero-shot scripts need lm-eval 0.3.0 and its C4 loader fails with new \texttt{datasets} versions. We will first check that we get the paper's dense perplexity of 5.68 on LLaMA-7B. One GitHub issue reports better numbers than the paper (20.54 instead of 24.55 for OWL with Wanda), while other users get 24.1 to 24.4. So we will report our own dense and uniform baselines next to the paper's numbers. We count a result as reproduced if the method ranking matches at every sparsity and the OWL gain is of similar size.

\section{Plans for Extension}

C2 was only tested on LLaMA-1 (7B to 65B), trained on at most 1.4T tokens. Modern small models see far more data, about 15T tokens for Llama~3 \citep{grattafiori2024llama3} and 18T for Qwen2.5 \citep{qwen2025qwen25}. Evidence for OWL on them is mixed. With SparseGPT at 70\%, OWL nearly halves Llama-3-8B perplexity, from 85.84 to 48.07 \citep{sieberling2025evopress}. But OWL slightly hurts a 1B Llama~3 model in a multilingual setting \citep{choenni2025mwanda}, and layerwise allocation gives only marginal gains on Llama-3-8B in \citet{jamshidi2026pals}. ATP \citep{huang2025atp} even argues for a schedule that simply increases with depth. As far as we know, no work tests OWL across sizes within one modern family or below 3B, and no work checks the LOD shape of these models. We ask two questions. Does OWL's benefit still grow as models shrink? Does the LOD predict when OWL helps?

\paragraph{E1: Scaling within a family.} We prune Qwen2.5 0.5B, 1.5B, 3B and 7B, and Llama-3.2-1B, Llama-3.2-3B and Llama-3.1-8B, at 50\%, 60\% and 70\%. Each model gets Wanda and SparseGPT with uniform and OWL allocation. We report WikiText-2 perplexity at every sparsity and mean zero-shot accuracy at 70\%, then plot the OWL gain against model size.

\paragraph{E2: Does the LOD predict the gain?} We compute the LOD of all ten dense models, including the three from the reproduction, and check for the U shape. We measure how uneven each LOD is with its coefficient of variation across blocks. A Spearman correlation then tests whether more uneven models gain more from OWL.

\paragraph{E3: Hyperparameter transfer.} $M$ and $\lambda$ were tuned for LLaMA-1, so a weak result could just mean bad settings. We run the paper's grid ($M \in \{3,5,7,10\}$, $\lambda \in \{0.02,0.05,0.08,0.1,0.2\}$) on Qwen2.5-1.5B and Llama-3.2-1B at 70\% with Wanda. We select on C4 validation perplexity, not WikiText-2 test, and rerun E1 at 70\% with the tuned values.

\paragraph{E4 (optional): Calibration language.} We calibrate Qwen2.5-1.5B and 7B on Vietnamese Wikipedia instead of English C4 (Wanda, uniform and OWL, 60\% and 70\%). We compare the LOD profiles and the perplexity on WikiText-2 and held-out Vietnamese text. Prior work on calibration language \citep{kurz2026limitations,choenni2025mwanda} did not study how OWL's allocation shifts.

\paragraph{Code changes.} The code calls each block in a way that breaks on \texttt{transformers} 4.48 and later. We record the arguments the model passes to each block and replay them, so it runs on current versions. We also use \texttt{AutoTokenizer} and run lm-eval in the same process, so no pruned checkpoint is saved.

\section{Resources}

No paid API is needed. Pruning needs direct access to the weights, and all models above are open-weight on Hugging Face (Llama~3 needs a free access request).

\begin{table}[t]
\centering
\small
\begin{tabular}{llrr}
\toprule
Exp. & Models & Runs & A100 h \\
\midrule
R1 & LLaMA-7B/13B, OPT-6.7B & 15 & 3 \\
R2 & LLaMA-7B & 28 & 2 \\
R3 & LLaMA-7B & 6 & 3 \\
R4, E2 & 10 dense models & 10 & 1 \\
E1 & 7 Qwen2.5 / Llama-3 & 84 & 12 \\
E3 & 2 small models & 40 & 2 \\
E4 & Qwen2.5-1.5B/7B & 16 & 2 \\
\midrule
Total & & 199 & 25 \\
\bottomrule
\end{tabular}
\caption{Planned runs (pruning plus evaluation) and estimated A100 hours.}
\label{tab:budget}
\end{table}

\paragraph{Memory.} A 7B model in fp16 takes 13.5\,GB. The calibration activations of one block add about 4\,GB, so the peak is about 18\,GB. The 7B and 8B models need a 24\,GB or larger GPU, LLaMA-13B needs about 31\,GB, and models up to 3B fit on a 16\,GB card.

\paragraph{Compute.} On an A100 the paper reports 266\,s for SparseGPT and under 1\,s for the Wanda metric on 7B, plus forward passes. Table~\ref{tab:budget} estimates 25 A100 hours, and we budget 50 to cover reruns. Our main resource is USF's GAIVI cluster (15 A100, 6 H100 and 10 L40S GPUs), which is free for students with a supervisor's approval. Its jobs can be preempted, but each run is short and saves its results at the end, so we lose at most one run.

\paragraph{Backup.} Kaggle gives two 16\,GB T4 GPUs for about 30 hours a week, enough for all models up to 3B and the E3 grid. Modal's free \$30 monthly credit buys about 14 A100 hours for the 7B SparseGPT runs. A rented RTX 4090 or A100 costs \$0.3 to \$1.8 per hour, so even a fully paid budget stays under \$70.

\bibliography{references}

\end{document}
